% Part: applied-modal-logics
% Chapter: epistemic-logic
% Section: truth-at-w

\documentclass[../../../include/open-logic-section]{subfiles}

\begin{document}

\olfileid{aml}{el}{trw}

\olsection{ఒక లోకం వద్ద సత్యం}

సాధారణ మోడల్ తర్కంలోలాగే, ప్రతి జ్ఞానసంబంధ నమూనా
దానిలోని ఏ లోకాల వద్ద ఏ \tetoken{సూత్రాలు}{!!{formula}s}
సత్యమో నిర్ణయిస్తుంది. సంబంధ అర్థవిజ్ఞానపు ప్రాథమిక
భావన కోసం ``నమూనా~$\mModel{M}$లో
\tetoken{సూత్రం}{!!{formula}}~$!A$ లోకం~$w$ వద్ద సత్యం''
అనే సంకేతాన్నే వాడతాం. ఆ సంబంధాన్ని ఆగమన పద్ధతిలో
నిర్వచిస్తాం; మోడల్ కాని కారకాలన్నింటికీ అది సాధారణ
మోడల్ తర్కంలోని నిర్వచనమే.

\begin{defn}\ollabel{defn:mmodels}
  నమూనా~$\mModel M$లో \emph{$w$ వద్ద
  \tetoken{ఒక సూత్రం}{!!a{formula}}~$!A$ సత్యం},
  సంకేతంగా $\mSat{M}{!A}[w]$, ఇలా ఆగమన పద్ధతిలో
  నిర్వచించబడుతుంది:
  \begin{enumerate}
  \tagitem{prvFalse}{\indcase{!A}{\lfalse}{$\mSat{M}{\lfalse}[w]$ ఎప్పుడూ కాదు}.}{}
  \tagitem{prvTrue}{\indcase{!A}{\ltrue}{$\mSat{M}{\ltrue}[w]$ ఎల్లప్పుడూ}.}{}
  \item $\mSat{M}{p}[w]$ అయితే, అప్పుడే $w \in V(p)$
  \tagitem{prvNot}{\indcase{!A}{\lnot !B}{$\mSat{M}{\indfrm}[w]$ అయితే, అప్పుడే
    $\mSat/{M}{!B}[w]$}.}{}
  \tagitem{prvAnd}{\indcase{!A}{(!B \land !C)}{$\mSat{M}{\indfrm}[w]$ అయితే, అప్పుడే
    $\mSat{M}{!B}[w]$ మరియు $\mSat{M}{!C}[w]$}.}{}
  \tagitem{prvOr}{\indcase{!A}{(!B \lor !C)}{$\mSat{M}{\indfrm}[w]$ అయితే, అప్పుడే
    $\mSat{M}{!B}[w]$ లేదా $\mSat{M}{!C}[w]$} (లేదా రెండూ).}{}
  \tagitem{prvIf}{\indcase{!A}{(!B \lif !C)}{$\mSat{M}{\indfrm}[w]$ అయితే, అప్పుడే
    $\mSat/{M}{!B}[w]$ లేదా $\mSat{M}{!C}[w]$}.}{}
  \tagitem{prvIff}{\indcase{!A}{(!B \liff !C)}{$\mSat{M}{\indfrm}[w]$ అయితే, అప్పుడే
    $\mSat{M}{!B}[w]$, $\mSat{M}{!C}[w]$ రెండూ సత్యం; లేదా
    $\mSat{M}{!B}[w]$, $\mSat{M}{!C}[w]$ రెండూ అసత్యం}.}{}
  \item\ollabel{defn:sub:mmodels-box}
    \indcase{!A}{\Knows_a !B}{$\mSat{M}{\indfrm}[w]$ అయితే, అప్పుడే
    $R_a ww'$ అయ్యే ప్రతి $w' \in W$కు $\mSat{M}{!B}[w']$}
  \end{enumerate}
\end{defn}

ఇక్కడే ప్రాప్యత సంబంధాలపై పరిమితులను ఆలోచించాలి.
ఎందుకంటే \olref{defn:sub:mmodels-box} నిబంధన ప్రకారం,
$R_a ww'$ అయ్యే~$w'$ ఏదీ లేకపోతే
\tetoken{సూత్రం}{!!a{formula}}~$\Knows_a !B$
లోకం~$w$ వద్ద సత్యమవుతుంది. సాధారణ మోడల్ తర్కంలోనూ
ఇదే నిబంధన: ఒక లోకానికి ఉత్తరవర్తులు లేకపోతే,
అక్కడ అన్ని $\Box$-సూత్రాలు శూన్యసత్యంగా నిలుస్తాయి.
$\Knows$ \emph{జ్ఞానాన్ని} సూచిస్తుందని భావిస్తే ఇది
చాలా ప్రతిసహజంగా కనిపిస్తుంది. ఒక కర్తకు ఒకే సమయంలో
$!A$, $\lnot !A$ రెండూ తెలిసే పరిస్థితులు
\emph{ఏవీ లేవు} అని సాధారణంగా భావిస్తాం.

ఒక పరిష్కారం: జ్ఞానసంబంధ తర్కంలో మన ప్రాప్యత సంబంధం
ఎల్లప్పుడూ \emph{స్వావర్తనమైనది} అయ్యేలా చూడటం.
వాస్తవ లోకం కర్త సమాచారానికి అనుకూలమే అన్న భావనకు
ఇది స్థూలంగా అనురూపం. నిజానికి జ్ఞానసంబంధ తర్కాలు
సాధారణంగా S5ను వాడతాయి; అయితే $\Knows_a$
సంబంధం ఖచ్చితంగా దేనిని సూచించాలనుకుంటున్నారనే దాన్ని
బట్టి ఇతరులు బలహీనమైన వ్యవస్థలను వాడవచ్చు.

\begin{figure}
  \begin{center}
    \begin{tikzpicture}[modal]
      \node[world] (w1) [label={[align=right]right:\mTrue{p}\\ \mFalse{q}}]{$w_1$};
      \node[world] (w2) [label={[align=right]right:\mTrue{p}\\ \mTrue{q}},
        above right=of w1]{$w_2$};
      \node[world] (w3) [label={[align=right]right:\mFalse{p}\\ \mFalse{q}},
        below right=of w1] {$w_3$};
      \draw[<->] (w1) to node {$a$} (w2);
      \draw[->,loop left] (w1) to node {$a,b$} (w1);
      \draw[<->] (w1) to [swap] node {$b$} (w3);
      \draw[->,loop above] (w2) to node {$a,b$} (w2) ;
      \draw[->,loop below] (w3) to node {$a,b$} (w3) ;
    \end{tikzpicture}
  \end{center}
  \caption{సరళమైన జ్ఞానసంబంధ నమూనా.}
  \ollabel{fig:simple}
\end{figure}

\begin{prob}
  కింది వాటిలో ఏవి
  \olref[aml][el][trw]{fig:simple}లో నిలుస్తాయో పరిశీలించండి:
  \begin{enumerate}
    \item $\mSat{M}{\lnot q}[w_1]$;
    \item $\mSat{M}{\Knows_a \lnot q}[w_1]$;
    \item $\mSat{M}{\Knows_b \lnot q}[w_1]$;
    \item $\mSat{M}{\Knows_b q \lor \Knows_b \lnot q}[w_2]$;
    \item $\mSat{M}{\Knows_a( \Knows_b q \lor \Knows_b \lnot q)}[w_2]$;
    \item $\mSat{M}{\EKnows_{\{a,b\}} \lnot q}[w_3]$;
    \end{enumerate}
\end{prob}

ఒక లోకం వద్ద సత్యానికి ప్రాథమిక నిర్వచనాన్ని
ఇచ్చాం. ఇక మోడల్ చెల్లుబాటు, అనుగమనం వంటి సాధారణ
మోడల్ తర్కంలోని ఇతర అర్థవిజ్ఞాన భావనలూ ఇక్కడికి
అలాగే వర్తిస్తాయి; మోడల్ కారకాల వ్యాఖ్యానాన్ని
కొత్తగా అర్థం చేసుకోవడమే తేడా.

ఇప్పుడు సామాన్య జ్ఞాన కారకం~$\CKnows_G$కు సత్య
షరతులను ఇవ్వవచ్చు. ఒక సంబంధం~$R$ యొక్క
\emph{సంక్రామక సంవృతం}~$R^+$ను
\olref[sfr][rel][ops]{sec} నుంచి గుర్తుచేసుకోండి:
\begin{align*}
  R^+ &= \bigcup_{ n \in \mathbb{N}} R^n,
\intertext{ఇక్కడ}
  R^0 & = R \text{ మరియు}\\
  R^{n+1} & = \Setabs{\tuple{x, z}}{\exists y (R^n xy \land Ryz)}.
\end{align*}
అప్పుడు, $G$ కర్తల సమూహమైతే,
అన్ని కర్తల ప్రాప్యత సంబంధాల సమ్మేళనపు
సంక్రామక సంవృతాన్ని $R_G = ( \bigcup_{b
\in G} R_b )^+$గా నిర్వచిస్తాం.

\begin{defn}
$G' \subseteq G$ అయితే, $R_{G'} w w'$ అయ్యే ప్రతి
$w'$కు $\mSat{M}{!A}[w']$ అయినప్పుడు, అప్పుడే
$\mSat{M}{\CKnows_{G'} !A}[ w]$ అని నిర్వచిస్తాం.
\end{defn}

\end{document}
